\documentclass[11pt]{article}
\usepackage[margin=1in]{geometry}
\usepackage{amsmath,amssymb,amsthm,mathtools}
\usepackage[T1]{fontenc}
\usepackage{lmodern,microtype,xurl}
\usepackage[colorlinks=true,linkcolor=blue,urlcolor=blue]{hyperref}
\newtheorem{theorem}{Theorem}[section]
\newtheorem{lemma}[theorem]{Lemma}
\newtheorem{proposition}[theorem]{Proposition}
\newtheorem{corollary}[theorem]{Corollary}
\theoremstyle{remark}\newtheorem{remark}[theorem]{Remark}
\newcommand{\E}{\mathbb E}
\newcommand{\Prb}{\mathbb P}
\newcommand{\R}{\mathbb R}
\newcommand{\norm}[1]{\left\lVert #1\right\rVert}
\title{Finite-width bounds for the globally margined tanh posterior}
\author{ChatGPT-6 Astra \and Kaarel H\"anni \and Dmitry Vaintrob}
\date{26 October 2026}
\begin{document}
\maketitle
\begin{abstract}
We give finite-width bounds for the actual Gaussian tanh posterior conditioned on
absolute margin greater than $1/2$ at every Boolean input and exact training
labels. A bounded-weight, globally margined width-$C$ teacher at depth $L$
is learned at every width $N\ge C$ with sample size essentially
$[Nd+LN^2]L\log(d+C+L+B+N)$, uniformly over canonical variance prefactors
$0<t\le1$. The proof combines a teacher neighborhood, a conditional energy
tail, and a margin-stable parameter net. For bounded-fan-in Boolean circuit
targets, a separate Gaussian mean-shift witness improves this to a
linear-width bound $O(N(d+CH)/t^2)$ above an explicit width threshold,
where $H=L-1$ and $C$ counts the circuit wires. This second bound retains
its inverse-scale dependence. Conversely, augmented-coordinate
symmetry forces $\Omega(\log d)$ samples for a constant target, at every
finite width and depth. We also exhibit a two-dimensional example in which
correct-label posterior probability strictly decreases when width increases
from one to two. These are finite-width statements. They neither remove
the learner-width dependence from the upper bound nor establish an
exponential-in-depth lower bound.
\end{abstract}

\section{Model and confidence criterion}

Let $X_d=\{-1,1\}^d$, $d\ge1$, with augmented input $(1,x)$.
There are $L\ge2$ matrices: $L-1$ tanh hidden layers, each of width $N$,
and a linear readout $z_W(x)$. There are no internal biases. Independent
Gaussian entries have variances
\[
 t^2/d,\qquad t^2/N,\qquad t^2/N^2
\]
in the first matrix, internal matrices, and readout, respectively.
The canonical range is $0<t\le1$. Write $\mu$ for this raw Gaussian law.
For a labeled sample $S=((x_i,f(x_i)))_{i=1}^n$, the posterior is
\[
 Q_S=\mu(\,\cdot\mid E_S),\qquad
 E_S=\{|z_W(x)|>1/2\ \forall x\in X_d\}
          \cap\{\operatorname{sign}z_W(x_i)=f(x_i)\ \forall i\}.
\]
This is exactly the result of conditioning first on global absolute
margin and then on the labels. Repeated examples add no new constraints.

For an input law $D$, put
\[
 p_S(x)=Q_S\{\operatorname{sign}z_W(x)=f(x)\},\qquad
 e(S)=\Prb_{x\sim D}\{p_S(x)\le2/3\}.
\]
The upper bound below controls $e(S)$ with high probability over an iid
sample. Its teacher hypothesis is a width-$C$ representation at the same
depth $L$, with physical weights bounded by $B\ge1$ and signed output
at least one everywhere. A margin-$1/2$ teacher is handled by doubling
its readout and replacing $B$ by $2B$.

\section{A sample upper bound quadratic in learner width}

\begin{theorem}\label{thm:upper}
There is a universal $c$ such that, for every teacher as above, every
$N\ge C$, every $0<t\le1$, every input law $D$, and
$0<\epsilon,\delta<1$, one has $\Prb_S\{e(S)\le\epsilon\}\ge1-\delta$
whenever
\[
 n\ge\frac c\epsilon\left[
 pL\log(d+C+L+B+N+\epsilon^{-1}+2)+\log(1/\delta)\right],
 \qquad p=N(d+2)+(L-2)N^2.
\]
The estimate is uniform in prefactor schedules and in limiting
strong-prior posteriors at fixed architecture.
\end{theorem}

\subsection{Energy coordinates and a stable teacher neighborhood}

Use coordinates
\[
 A_1=\sqrt{d/N}\,W_1,\qquad A_j=W_j\ (2\le j<L),
 \qquad A_L=\sqrt N\,v.
\]
Their $p$ entries are independent $N(0,\sigma^2)$ with
$\sigma^2=t^2/N\le1/N$. Write
$\kappa(A)=\sum_j\norm{A_j}_F^2=\norm A_2^2$.
Equivalently, the network takes input $u_0=(1,x)/\sqrt d$ and applies
each hidden matrix followed by
\[
 \phi_N(s)=\tanh(\sqrt N s)/\sqrt N.
\]
This activation fixes zero and is 1-Lipschitz, and
$\norm{u_0}\le\sqrt2$. The last layer remains linear.

\begin{lemma}\label{lem:teacher}
The teacher has an exactly equivalent width-$N$ realization $A_*$ with
\[
 \kappa(A_*)\le K_0:=2B^2\{d(d+1)+(L-1)C^2\}.
\]
On the energy ball $\kappa(A)\le R$, the output at every input is
Lipschitz in the Euclidean parameter norm, with constant
\[
 \Lambda_R=\sqrt{2L}\max\{1,R/(L-1)\}^{(L-1)/2}.
\]
\end{lemma}
\begin{proof}
Pad the teacher to width $C$, and replicate each hidden neuron either
$\lfloor N/C\rfloor$ or $\lceil N/C\rceil$ times. Retain its incoming
weights and divide outgoing weights by its replication count. Outputs
are preserved exactly. The weighted first and last energies increase
by at most two; each internal squared Frobenius norm increases by at
most the largest ratio of two replication counts, also at most two.
At width $C$, the three energy contributions are at most
$d(d+1)B^2$, $(L-2)C^2B^2$, and $C^2B^2$, proving the first claim.

The derivative with respect to matrix block $j$ has Frobenius norm at
most $\sqrt2$ times the product of the other $L-1$ block norms.
Arithmetic--geometric mean bounds that product by
$(R/(L-1))^{(L-1)/2}$. Sum the squared block derivative bounds and
integrate along the segment joining two parameters. The energy ball
is convex, so the same bound holds on the entire segment.
\end{proof}

Set $K=2K_0+2$ and $r=1/(4\Lambda_K)\le1$. The Euclidean $r$-ball
around $A_*$ lies in the energy-$K$ ball, and every parameter in it
has signed output at least $3/4$ at every input. Thus it belongs to
$E_S$ for every sample labeled by $f$.

\subsection{A conditional energy tail at every width}

Inside this teacher neighborhood take the coordinate cube with half-side
$h=\min\{r/\sqrt p,\sigma\}$. Its Gaussian density is at least
$(2\pi\sigma^2)^{-p/2}e^{-K/(2\sigma^2)}$, so
\begin{equation}\label{eq:denominator}
 \mu(E_S)\ge e^{-K/(2\sigma^2)-pT},\qquad
 T=\log\!\left[\sqrt{\pi/2}
                   \max\{1,\sqrt{p/N}/r\}\right].
\end{equation}
Indeed its volume is $(2h)^p$, and $\sigma\sqrt p\le\sqrt{p/N}$.
Clipping the cube at $\sigma$ makes $T$ independent of the prefactor;
there is no $\log(1/t)$ loss.

The elementary Gaussian Chernoff bound gives
\[
 \mu\{\kappa(A)>R\}\le2^{p/2}e^{-R/(4\sigma^2)}.
\]
For $0<\beta<1$, define
\begin{equation}\label{eq:energy-radius}
 R_\beta=2K+\frac4N
       \left[p\{T+\tfrac12\log2\}+\log(1/\beta)\right].
\end{equation}
Divide the tail bound by \eqref{eq:denominator}, and use
$\sigma^2\le1/N$, to obtain the sample-uniform posterior estimate
\begin{equation}\label{eq:posterior-tail}
 Q_S\{\kappa(A)>R_\beta\}\le\beta.
\end{equation}
Both conditioning events remain present in this calculation.

Writing $s=d+C+L+B+2$, the definitions imply
\[
 T\le c[L\log s+\log(d+LN+2)],
\]
\[
 R_\beta\le c\left[K+(d+LN)\{L\log s+\log(d+LN+2)\}
                    +N^{-1}\log(1/\beta)\right].
\]
In particular this radius is polynomial in the ordinary resources and
$N$, uniformly over all canonical scales.

\subsection{Margin-stable signs and posterior confidence}

Cover the energy-$R_\beta$ ball by a Euclidean net of radius
$1/(8\Lambda_{R_\beta})$, with centers in the ball. The number of
centers satisfies
\[
 M\le(1+16\sqrt{R_\beta}\Lambda_{R_\beta})^p,
 \qquad
 \log M\le cpL\log(s+N+\beta^{-1}).
\]
Every globally clean-margin network in this ball has the same signs
on the entire cube as a net center: the uniform output change is at
most $1/8$. These centers therefore define a fixed class of at most
$M$ truth tables, before the sample or input law is chosen. Centers
with a zero output may use any fixed sign convention; those selected
for a clean-margin network have no such zero.

A finite-class union bound shows that, with sample probability at least
$1-\delta$, every center fitting the labels has population error at
most $\alpha$, provided $n\ge[\log M+\log(1/\delta)]/\alpha$.
On this event all low-energy posterior networks have error at most
$\alpha$. Equation~\eqref{eq:posterior-tail} bounds the remaining
posterior mass by $\beta$, so posterior Gibbs error is at most
$\alpha+\beta$. Every point counted by $e(S)$ contributes at least
$1/3$ to Gibbs error. Thus $e(S)\le3(\alpha+\beta)$.
Taking $\alpha=\beta=\epsilon/6$ proves Theorem~\ref{thm:upper}.

The sample event and the Gibbs bound hold uniformly in $t$. On the
finite truth-table simplex, Gibbs error passes continuously to any
limiting strong-prior posterior. Applying the same final inequality
proves the stated limiting version without asserting continuity of
the confidence-threshold event itself.

\begin{corollary}\label{cor:exponential-width}
If $d,C,B$ are polynomially bounded in $L$, and a fixed positive
accuracy/confidence requirement needs at least $e^{cL}$ samples along
a canonical finite-width sequence, then
\[
 \log N\ge cL/2-O(\log L).
\]
\end{corollary}
\begin{proof}
If $\log N>cL$ there is nothing to prove. Otherwise the sufficient
sample bound in Theorem~\ref{thm:upper} is $N^2$ times a polynomial
in $L$. Rearranging gives the claim.
\end{proof}

\section{A linear-width bound for Boolean circuit targets}
\label{sec:boolean-mean-shift}

For targets specified by digital gates, a direct Gaussian mean shift
improves the preceding width dependence. This theorem retains an
explicit $t^{-2}$ factor; it is not a uniform cold-limit bound.
Write $H=L-1$. Suppose $f$ has a layered AND/OR/NOT/COPY circuit
with at most $H$ stages and width at most $C$, including a constant
$+1$ wire. Pad shorter circuits by COPY gates.

\begin{theorem}\label{thm:boolean-linear-width}
If
\[
 N\ge1600C\{d\log2+\log(8CH)\},
\]
then the raw Gaussian prior satisfies
\[
 \mu\{f(x)z_W(x)>1/2\ \forall x\in X_d\}
 \ge \exp[-2002N(d+CH)/t^2].
\]
For any common training and test input law, the actual posterior
therefore has confidence-failure mass at most $\epsilon$ with
sample probability at least $1-\delta$ whenever
\[
 n\ge\frac6\epsilon\left[
       \frac{2002N(d+CH)}{t^2}
       +\log\frac6{\epsilon\delta}\right].
\]
\end{theorem}

\begin{proof}
Partition each hidden layer into $C$ groups whose sizes $n_g$
are $\lfloor N/C\rfloor$ or $\lceil N/C\rceil$, hence
$n_g\ge N/(2C)$. A group represents one circuit wire by its
average activation. Use gain $a=10$ and the affine gate formulas
\[
 u+v-1,\quad u+v+1,\quad -u,\quad u
\]
for AND, OR, NOT, and COPY. A constant term is supplied by the
actual previous constant group, or by the augmented input at the
first layer. Coincident inputs may be combined; the coefficient
absolute sum remains at most three.

Define a product Gaussian comparison law $Q$ with the original
variances and deterministic shifted means. A first-layer gate row
has mean $a$ times its affine formula's coefficients. At later
layers, put mean $ac_j/n_j$ on every coordinate of predecessor
group $j$, where $c_j$ is that formula's coefficient. All other
weight means are zero. The constant group copies its predecessor
constant, starting from the augmented constant input. Set the
readout mean to $1/n_{\rm out}$ on the output group and zero
elsewhere. All background random weights remain present.

Suppose the predecessor group averages have the correct signs
and signed magnitudes at least $0.9$. Since all tanh values lie
in $[-1,1]$, AND and OR retain signed affine margin at least
$2(0.9)-1=0.8$, including error in the constant group. COPY and
NOT retain margin at least $0.9$. The first-layer margin is at
least one. Thus each controlled row has signed preactivation
mean at least eight. Conditional on previous layers its variance
is at most two initially and at most one later, uniformly for
$t\le1$. For centered Gaussian $Z$ of variance at most two,
\[
 \E\tanh(8+Z)\ge1-2e^{-4}-2e^{-8}>0.95.
\]
Indeed $\Prb(Z<-4)\le e^{-4}$ and $\tanh(8+Z)\ge1-2e^{-8}$
on the complementary event. Symmetry treats a negative gate bit.

Conditional Hoeffding concentration gives probability at most
$e^{-n_g/800}$ that one signed group average is below $0.9$ at
one input. Sequential union bounds over all layers, groups, and
cube inputs give hidden propagation failure at most
\[
 CH2^d e^{-N/(1600C)}\le1/8.
\]
No independence between distinct cube inputs is used. Conditional
on the hidden layers, centered readout noise has variance at most
$t^2/N$ at each input, so
\[
 Q\{\max_x|z(x)-\E[z(x)\mid\text{hidden layers}]|>1/4\}
 \le2^{d+1}e^{-N/(32t^2)}\le1/8.
\]
On the joint good event every signed output is at least
$0.9-0.25=0.65$. Thus $Q(E_f)\ge3/4$ for the globally correct
margin event $E_f$ in the theorem.

We charge every mean shift. The first-layer relative entropy is
at most $450dN/t^2$. An internal row has squared mean norm at most
\[
 a^2\sum_jc_j^2/n_j\le18a^2C/N,
\]
so its entropy is at most $900C/t^2$. The readout contributes
at most $NC/t^2$. Hence
\[
 \mathrm{KL}(Q\Vert\mu)
 \le\frac N{t^2}\{450d+900C(H-1)+C\}
 \le\frac{1000N(d+CH)}{t^2}.
\]
For any event $E$, binary relative entropy gives
\[
 \mathrm{KL}(Q\Vert\mu)
 \ge Q(E)\log(1/\mu(E))-\log2.
\]
Using $Q(E_f)\ge1/2$ proves the raw target-mass bound with
constant $2002$.

Let $\pi$ be the prior already conditioned on global absolute
margin. Since $E_f$ is contained in that event, its truth-table
mass under $\pi$ is at least
$p_*:=\exp[-2002N(d+CH)/t^2]$. This is a lower bound for every
target-labeled sample's version-space normalizer. The expected
$\pi$-mass of classifiers fitting an iid sample and having
population error greater than $\epsilon/6$ is at most
$e^{-n\epsilon/6}$. Markov therefore gives
\[
 \Prb_S\{Q_S(\operatorname{err}>\epsilon/6)>\epsilon/6\}
 \le\frac{6e^{-n\epsilon/6}}{\epsilon p_*}\le\delta.
\]
Outside this exceptional sample event the posterior Gibbs error
is at most $\epsilon/3$. A point with correct-label probability
at most $2/3$ contributes at least $1/3$ to that Gibbs error, so
its total probability mass is at most $\epsilon$. This argument
uses the actual posterior mixture and both conditioning events.
\end{proof}

\begin{remark}[Layerwise resources and exponent scaling]
If the circuit widths are $C_1,\ldots,C_H$, put
$C_{\max}=\max_jC_j$ and $G=\sum_jC_j$, including all constant
and padding COPY wires. Separate balanced partitions at each
layer give internal entropy at most $900NC_{j-1}/t^2$ at stage
$j$, and readout entropy at most $NC_H/t^2$. Thus replace $CH$
by $G$ in the theorem and its width threshold by
\[
 1600C_{\max}\{d\log2+\log(8G)\}.
\]
Under the common exponent convention
$t=N^{(1-\gamma)/2}$, the width factor $N/t^2$ equals $N^\gamma$.
One may take the smaller of this bound and Theorem~\ref{thm:upper}.
At $\gamma=1$ the new dependence is linear in $N$ above the
displayed threshold. For $C_{\max}\le N$ below that threshold,
Theorem~\ref{thm:upper} and the bounded-weight digital realization
give a coarser $N$ times polynomial-circuit-resource bound, since
the extra factor $N$ is then polynomially bounded. This comparison
does not assert the sharper coefficient $d+G$ below the threshold.
\end{remark}

The polynomial-depth statement here concerns a digital Boolean
implementation at the available learner depth. A tanh teacher
with soft intermediate states need not have such an implementation
at comparable resources. Nor does the explicit $t^{-2}$ factor
permit taking a cold limit in this particular bound. Exponentially
large learner widths can still make its sample bound exponential
in depth.

\section{An exact symmetry lower bound for a constant target}

Here $f\equiv+1$ and $D$ is uniform on the cube. The argument works
for every positive prefactor $t$, every finite width and depth, and
more generally whenever the first-layer law is invariant under signed
permutations of its augmented input columns and is independent of the
remaining layers. Only odd hidden activations and a linear readout
without internal biases are required, with a conditioning event of
positive probability.

For a confidence threshold $c>1/2$, write
$e_c(S)=\Prb_x\{Q_S(\operatorname{sign}z_W(x)=+1)\le c\}$.
A sample coordinate column is \emph{constant} if all its observed
values are the same sign.

\begin{theorem}\label{thm:symmetry}
If a nonempty sample has $k$ constant coordinate columns, then
\[
 e_c(S)\ge\frac12-2^{-k-1}\binom{k}{\lfloor k/2\rfloor}
 \qquad(c>1/2).
\]
For $n\ge1$ uniform iid examples, their number satisfies
$K\sim\operatorname{Binomial}(d,2^{1-n})$. In particular
\[
 \Prb_S\{e_c(S)\ge1/4\}\ge1-(1-2^{1-n})^d.
\]
For fixed $0<\epsilon<1/2$ and $0<\delta<1$, obtaining
$e_c(S)\le\epsilon$ with probability at least $1-\delta$ requires
$n\ge\log_2d-O_{\epsilon,\delta}(1)$.
\end{theorem}
\begin{proof}
First suppose coordinate $j$ has observed value $s\in\{-1,1\}$.
Let $A$ exchange augmented coordinates $0,j$ and multiply both by $s$.
Define a cube involution $T$ by
\[
 (Tx)_j=x_j,\qquad (Tx)_\ell=sx_jx_\ell\quad(\ell\ne j).
\]
Then $A(1,x)=sx_j(1,Tx)$. The whole network is odd in its augmented
input, so applying $A$ to the first-layer columns transforms its output
to $z_A(x)=sx_jz(Tx)$. This preserves the raw prior and the global
absolute-margin event. At every observed point $x_j=s$, it also has
$Tx=x$ and $sx_j=1$, preserving every training constraint.
Thus it preserves the exact posterior, including all its parameter
components. If $p(x)=Q_S\{z_W(x)>0\}$, then on $x_j=-s$,
\[
 p(Tx)=1-p(x).
\]
The involution preserves uniform measure on that half cube. Since
$c>1/2$, at most half of it can have $p>c$, giving $e_c(S)\ge1/4$.

For the full bound, flip the $k$ constant columns so all observed
values become $+1$. This preserves the uniform input law and the
constant target. Permutations of these columns and augmented column
$0$ preserve every training constraint. Permuting the $k$ ordinary
columns shows that the posterior vote depends on them only through
their number $r$ of negative signs, and on the outside coordinates
$y$. Call it $p_r(y)$, and let $w_r=2^{-k}\binom kr$.
For $1\le r\le k$, exchange column $0$ with one negative coordinate
and use oddness to renormalize the augmented first coordinate. This gives
\begin{equation}\label{eq:paired-votes}
 p_{k+1-r}(-y)=1-p_r(y).
\end{equation}
For two distinct paired levels, their total failure mass is at least
$\min(w_r,w_{k+1-r})=\min(w_r,w_{r-1})$. At a self-paired middle
level, \eqref{eq:paired-votes} gives failure on at least half the
level; if no outside coordinates remain, it forces $p_r=1/2$ and
only strengthens this estimate. Counting every pair twice yields
\[
 e_c(S)\ge\frac12\sum_{r=1}^k\min(w_r,w_{r-1})
         =\frac12(1-\max_r w_r).
\]
The last identity follows from unimodality of the binomial weights.
For $k=0$ the stated bound is zero.

Under uniform iid sampling the $d$ coordinate columns are independent,
and each is constant with probability $2^{1-n}$. This proves the
binomial and single-column assertions. For the final statement choose
a fixed $k_0$ making the displayed error bound greater than $\epsilon$.
If $d2^{1-n}$ exceeds a sufficiently large constant depending on
$k_0,\delta$, a binomial Chernoff bound makes $K\ge k_0$ with
probability greater than $\delta$. Rearranging proves the necessity.
\end{proof}

A stronger asymptotic form follows directly. If
$n\le(1-\eta)\log_2d$ for fixed $\eta>0$, then with probability at
least $1-e^{-d^\eta/4}$ one has $K\ge d^\eta$, and hence
\[
 e_c(S)\ge1/2-O(d^{-\eta/2}).
\]
The constant target has a bounded-weight realization of width at most
two at every depth: use the augmented constant input, carry a positive
scalar with weight two through each hidden layer, and use readout
weight two. Its activations stay above $1/2$ and its output exceeds one.
Thus the lower bound already applies to a simple globally margined teacher.

\begin{remark}[Scope of the symmetry argument]
This concerns correct-label confidence, not an unspecified tie convention
for ordinary majority prediction. Its lower bound is logarithmic in
dimension. The group of signed augmented-coordinate permutations has
at most $2^{d+1}(d+1)!$ elements. For any target, its orbit under these
input-label actions is therefore a concept class of that size, and
ordinary realizable consistency controls this orbit after
$O((d\log(d+2)+\log(1/\delta))/\epsilon)$ samples. This does not
bound the entire neural posterior. It explains why indistinguishability
within this exact symmetry orbit alone cannot yield exponential-depth
hardness when $d$ is polynomial in depth.
\end{remark}

\section{A caution against pointwise monotonicity in width}

\begin{proposition}\label{prop:nonmonotone}
For the actual Gaussian posterior at positive finite scales, correct-label
posterior probability at an unseen input need not be nondecreasing with
learner width. No finite-factor domination of larger-width wrong-label
probability by smaller-width wrong-label probability holds uniformly.
\end{proposition}
\begin{proof}
Take $d=2$, $L=2$, and target
$f(x)=\operatorname{sign}(1+x_1+x_2)$. The width-one teacher
$z_f(x)=\tanh(1+x_1+x_2)$ has all weights one and strict global
margin greater than $1/2$. Observe the three labels
\[
 (1,-1)\mapsto+1,\qquad(-1,1)\mapsto+1,\qquad(-1,-1)\mapsto-1.
\]
The only completions are $f$ and $g(x)=-x_1x_2$. A width-one
one-hidden-layer tanh network has an affine-threshold sign function,
whereas $g$ is not an affine threshold. Consequently the width-one
posterior assigns probability one to the correct label at $(1,1)$.
Its conditioning event has positive probability by a strict-margin
neighborhood of the displayed teacher.

At width two, the network
\[
 z_g(x)=3\tanh(10(1+x_1+x_2))
       -6\tanh(0.4+0.2x_1+0.2x_2)
\]
represents $g$ with strict global margin greater than $1/2$.
For $s=x_1+x_2=-2,0,2$, its values are respectively
\[
 -3\tanh10,\qquad3\tanh10-6\tanh0.4,\qquad
 3\tanh30-6\tanh0.8.
\]
The first is below $-2.97$, the second exceeds $0.57$, and the
third is below $-97/125$. These inequalities follow from
$\tanh10>0.99$, $\tanh u\le u$, and
$\tanh u\ge u-u^3/3$ for $u\ge0$.
The function $f$ also has a strict-margin width-two realization by
zero padding. Both therefore have open parameter neighborhoods of
positive Gaussian mass at every finite positive scale. After global
margin and sample conditioning, both completions retain positive mass.
The width-two correct-label probability at $(1,1)$ is strictly between
zero and one, proving both assertions.
\end{proof}

This example does not show that the width-two correct-label probability
crosses below $2/3$, and gives no uniform lower bound on its wrong-label
mass or strictness after a cold limit. A specialized comparison for
confidence-failure risk is not ruled out. At one hidden layer and fixed
$t$, the raw width-$kN$ output is an average of $k$ independent raw
width-$N$ outputs. That averaging identity does not survive as a
support-preserving comparison of globally margined sign posteriors.

\section{What the bounds leave open}

The general upper theorem and the stronger Boolean-circuit upper bound
supply actual-posterior obstructions to depth hardness at small widths; the lower theorem supplies a dimension-dependent
confidence obstruction already for a constant target. They are not
matching bounds. The general upper bound retains quadratic learner-width dependence;
the digital-circuit bound is linear in width at the critical scale, and the symmetry lower gives only $\Omega(\log L)$ when
$d$ is polynomial in $L$. The width counterexample blocks one simple
interpolation principle without proving a confidence lower bound.
The original question of a universal polynomial sample bound versus an
exponential-depth counterexample for canonical tanh remains open.
\end{document}
